\documentclass[tikz,border=3mm,multi=tikzpicture]{standalone}
\usepackage{eleve-matematica-ef1}

\begin{document}

% FIGURA 1 — fig-01-decimos-e-centesimos
\begin{tikzpicture}[matematica figura, x=0.72cm, y=0.72cm]
  \path[use as bounding box] (-0.20,-0.30) rectangle (10.60,12.10);
  \fill[matemalaranjaclaro] (0,8.00) rectangle (1,10.00);
  \draw[matematica contorno] (0,8.00) rectangle (10,10.00);
  \foreach \x in {1,...,9}\draw[matematica divisao] (\x,8.00)--(\x,10.00);
  \node[matematica valor] at (5,10.85) {$1$ de $10$ partes};
  \node[matematica resultado] at (5,7.15) {$\frac{1}{10}$};

  \foreach \x in {0,...,9}
    \foreach \y in {0,...,9}{
      \pgfmathtruncatemacro{\n}{10*\y+\x}
      \ifnum\n<20\fill[matemazulclaro] (\x*0.72,0.20+\y*0.55) rectangle ++(0.72,0.55);\fi
    }
  \draw[matematica contorno] (0,0.20) rectangle (7.20,5.70);
  \foreach \x in {1,...,9}\draw[matematica divisao] (\x*0.72,0.20)--(\x*0.72,5.70);
  \foreach \y in {1,...,9}\draw[matematica divisao] (0,0.20+\y*0.55)--(7.20,0.20+\y*0.55);
  \node[matematica resultado] at (8.70,2.95) {$\frac{20}{100}$};
\end{tikzpicture}

% FIGURA 2 — fig-02-fracao-e-escrita-decimal
\begin{tikzpicture}[matematica figura, x=1cm, y=1cm]
  \path[use as bounding box] (-0.15,-0.20) rectangle (9.60,7.80);
  \foreach \i in {0,...,9}\draw[matematica contorno] ({0.15+0.70*\i},5.05) rectangle ++(0.70,1.30);
  \fill[matemalaranjaclaro] (0.19,5.09) rectangle (0.81,6.31);
  \draw[matematica contorno] (0.15,5.05) rectangle (7.15,6.35);
  \foreach \x in {0.85,1.55,2.25,2.95,3.65,4.35,5.05,5.75,6.45}
    \draw[matematica divisao] (\x,5.05)--(\x,6.35);
  \node[matematica valor] at (8.35,5.70) {$\frac{1}{10}=0{,}1$};

  \foreach \x in {0,...,9}
    \foreach \y in {0,...,9}
      \draw[matematica divisao] ({0.15+0.55*\x},{0.40+0.42*\y}) rectangle ++(0.55,0.42);
  \fill[matemazulclaro] (0.19,0.44) rectangle (0.66,0.78);
  \draw[matematica contorno] (0.15,0.40) rectangle (5.65,4.60);
  \node[matematica valor] at (7.55,2.50) {$\frac{1}{100}=0{,}01$};
\end{tikzpicture}

% FIGURA 3 — fig-03-zeros-na-escrita-decimal
\begin{tikzpicture}[matematica figura, x=1cm, y=1cm]
  \path[use as bounding box] (-0.15,-0.20) rectangle (10.80,6.30);
  \foreach \i in {0,...,9}{
    \ifnum\i<5\fill[matemalaranjaclaro] (0.10+0.42*\i,1.10) rectangle ++(0.42,4.20);\fi
    \draw[matematica divisao] (0.10+0.42*\i,1.10) rectangle ++(0.42,4.20);
  }
  \draw[matematica contorno] (0.10,1.10) rectangle (4.30,5.30);
  \node[matematica valor] at (2.20,0.45) {$0{,}5$};

  \foreach \x in {0,...,9}
    \foreach \y in {0,...,9}{
      \pgfmathtruncatemacro{\n}{10*\y+\x}
      \ifnum\n<50\fill[matemazulclaro] (6.00+0.42*\x,1.10+0.42*\y) rectangle ++(0.42,0.42);\fi
      \draw[matematica divisao] (6.00+0.42*\x,1.10+0.42*\y) rectangle ++(0.42,0.42);
    }
  \draw[matematica contorno] (6.00,1.10) rectangle (10.20,5.30);
  \node[matematica valor] at (8.10,0.45) {$0{,}50$};
  \node[matematica resultado] at (5.15,3.20) {$=$};
\end{tikzpicture}

\end{document}
